%@ pages 377-379
%@ status translated
% Chapter III, section 4: Categories with generators and exact inductive
% limits (pp. 377--379).  Transcribed from the page images in
% reference/source.pdf, not from the OCR text layer.  The section ends
% at the top of p. 379 (Proposition 10 and the paragraph on simples);
% section 5 begins mid-p. 379.
%
% Numbering is automatic (src/preamble.tex).  This section prints
% Proposition 8, Proposition 9 and Proposition 10 (the chapter's run of
% propositions reaches Proposition 7 on p. 377), Corollary 1 and
% Corollary 2 (the corollary counter restarts at every section) and one
% corollary which Gabriel prints without a number, hence
% `unnumberedcorollary'.  No lemma is stated in this section.
%
% No emendations.

\section{Categories with Generators and Exact Inductive Limits}

--- We shall use the notations of the preceding paragraph.

\begin{proposition}
Let $\mathbf{C}$ be a thick subcategory of a $\Univ$-category
$\mathbf{A}$ abelian with generators and exact inductive limits. The
following assertions are equivalent:
\begin{description}
\item[\normalfont a.] $\mathbf{C}$ is a localizing subcategory of
      $\mathbf{A}$.
\item[\normalfont b.] The inductive limit (in $\mathbf{A}$) of an
      inductive system of objects of $\mathbf{C}$ belongs to
      $\mathbf{C}$.
\end{description}
\end{proposition}

$(a)\Rightarrow(b)$: let indeed $(M_i,\varphi_{ji})$ be an inductive
system of objects of $\mathbf{C}$, the set of indices belonging to the
universe $\Univ$. Let $M$ be the inductive limit of this system and
$\varphi_i$ the canonical morphism from $M_i$ into $M$. By
Corollary~1 to Proposition~6, $M$ contains a subobject $M'$ which
belongs to $\mathbf{C}$ and which contains every subobject of $M$
belonging to $\mathbf{C}$; in particular, $M'$ contains
$\operatorname{Im}\varphi_i$ for every $i$, and is therefore equal to
$M$.

$(b)\Rightarrow(a)$: let indeed $M$ be an object of $\mathbf{A}$. If
$P$ and $N$ are two subobjects of $M$ belonging to $\mathbf{C}$, then
$P+N$ is isomorphic to a quotient of $P\oplus N$. It follows that
$P+N$ belongs to $\mathbf{C}$ and that the subobjects of $M$
belonging to $\mathbf{C}$ form a filtered increasing family of
subobjects of $M$. By $(b)$, the supremum of these subobjects belongs
to $\mathbf{C}$ and is the largest subobject of $M$ belonging to
$\mathbf{C}$; whence the result, thanks to Corollary~1 to
Proposition~6.

\begin{unnumberedcorollary}
Let $\mathbf{A}$ be a $\Univ$-category abelian with generators and
exact inductive limits, and $(I_j)$ a family of injective objects of
$\mathbf{A}$. The objects $M$ of $\mathbf{A}$ such that
$\operatorname{Hom}_{\mathbf{A}}(M,I_j)$ is zero for every $j$ are the
objects of a localizing subcategory of $\mathbf{A}$. Conversely, every
localizing subcategory of $\mathbf{A}$ can be defined in this way.
\end{unnumberedcorollary}

Let indeed $\mathbf{C}$ be the full subcategory of $\mathbf{A}$ whose
objects $M$ are such that $\operatorname{Hom}_{\mathbf{A}}(M,I_j)$ is
zero for every $j$. If $0\to M'\to M\to M''\to0$ is an exact
sequence of $\mathbf{A}$, we also have the exact sequences
\[
0\to\operatorname{Hom}_{\mathbf{A}}(M'',I_j)
\to\operatorname{Hom}_{\mathbf{A}}(M,I_j)
\to\operatorname{Hom}_{\mathbf{A}}(M',I_j)\to0.
\]
These exact sequences show that $M$ belongs to $\mathbf{C}$ if and
only if $M'$ and $M''$ belong to $\mathbf{C}$. In other words,
$\mathbf{C}$ is a thick subcategory of $\mathbf{A}$.

On the other hand, every inductive system $(M_i,\varphi_{ji})$ of
$\mathbf{A}$ gives rise to the ``equalities''
\[
\operatorname{Hom}_{\mathbf{A}}\!\left(\varinjlim_i M_i,I_j\right)
\cong \varprojlim_i\operatorname{Hom}_{\mathbf{A}}(M_i,I_j).
\]
It follows that $\varinjlim_i M_i$ belongs to $\mathbf{C}$ if the
$M_i$ belong to $\mathbf{C}$.

Conversely, if $\mathbf{C}$ is a localizing subcategory of
$\mathbf{A}$, then $\mathbf{A}/\mathbf{C}$ is a category with
injective envelopes; an object $M$ of $\mathbf{A}$ therefore belongs
to $\mathbf{C}$ if and only if $\operatorname{Hom}_{\mathbf{A}}(M,I)$
is zero for every injective $\mathbf{C}$-closed object $I$.

\begin{proposition}
Let $\mathbf{C}$ be a localizing subcategory of a $\Univ$-category
$\mathbf{A}$ abelian with generators and exact inductive limits. The
categories $\mathbf{C}$ and $\mathbf{A}/\mathbf{C}$ are
$\Univ$-categories with generators and exact inductive limits. Moreover,
the canonical functor $T$ from $\mathbf{A}$ into
$\mathbf{A}/\mathbf{C}$ commutes with inductive limits.
\end{proposition}

Let $(U_i)_{i\in I}$ be a family of generators of $\mathbf{A}$. The
quotients, belonging to $\mathbf{C}$, of the objects $U_i$ form a
family of generators of $\mathbf{C}$. Moreover, every inductive system
of objects of $\mathbf{C}$ has an inductive limit in $\mathbf{A}$.
Proposition~8 shows that this inductive limit still belongs to
$\mathbf{C}$. It follows that $\mathbf{C}$ is a $\Univ$-category with
generators and exact inductive limits.

Lemma~5 (\S~2) proves that $\mathbf{A}/\mathbf{C}$ is a
$\Univ$-category with generators. Consider now an inductive system of
objects of $\mathbf{A}$, say $(M_i,\varphi_{ji})_{i\in I}$. Let
$\varphi_i$ be the canonical morphism from $M_i$ into
$\varinjlim_i M_i$. We then have the ``equalities'':
\[
\operatorname{Hom}_{\mathbf{A}/\mathbf{C}}\!\left(T\varinjlim_i M_i,N\right)
= \operatorname{Hom}_{\mathbf{A}}\!\left(\varinjlim_i M_i,SN\right)
= \varprojlim_i\operatorname{Hom}_{\mathbf{A}}(M_i,SN)
= \varprojlim_i
   \operatorname{Hom}_{\mathbf{A}/\mathbf{C}}(TM_i,N).
\]
These equalities show that $(T\varinjlim_i M_i,T\varphi_i)$ is an
inductive limit of the inductive system $(TM_i,T\varphi_i)$. The
functor $T$ therefore commutes with inductive limits.

Let now $(N_i,\psi_{ji})_{i\in I}$ be an inductive system of objects
of $\mathbf{A}/\mathbf{C}$ and let $\psi_i$ be the canonical morphism
from $SN_i$ into $\varinjlim_i SN_i$. By the above, the object
$T\varinjlim_i SN_i$ (together with the morphisms
$(T\psi_i)\circ\Phi(N_i)^{-1}$) is an
inductive limit of $(N_i,\psi_{ji})$. We shall write briefly
\[
\varinjlim_i N_i = T\varinjlim_i SN_i.
\]

This last ``formula'' shows in particular that an inductive limit
of monomorphisms is a monomorphism; this ends the proof.

\begin{corollary}
Let $\mathbf{C}$ be a localizing subcategory of a locally noetherian
$\Univ$-category. Then $\mathbf{C}$ and $\mathbf{A}/\mathbf{C}$ are
locally noetherian $\Univ$-categories. Moreover, the section functor
commutes with inductive limits.
\end{corollary}

It remains to show that $S$ commutes with inductive limits. If
$(N_i,\psi_{ji})$ is an inductive system of objects of
$\mathbf{A}/\mathbf{C}$, $\varinjlim_i N_i$ is equal to
$T\varinjlim_i SN_i$. It therefore suffices to prove that the
functorial morphism $\Psi$ defines an isomorphism from
$\varinjlim_i SN_i$ onto $ST\varinjlim_i SN_i$: this amounts to saying
that $\varinjlim_i SN_i$ is $\mathbf{C}$-closed.

This last assertion follows from Lemma~1~(a), \S2: with the notation
of that lemma, it indeed suffices to verify condition~(a) when $P$
and $Q$ are noetherian objects. The groups
$\operatorname{Hom}_{\mathbf{A}}(P,\varinjlim_i SN_i)$ and
$\operatorname{Hom}_{\mathbf{A}}(Q,\varinjlim_i SN_i)$ are then
identified respectively with
$\varinjlim_i\operatorname{Hom}_{\mathbf{A}}(P,SN_i)$ and
$\varinjlim_i\operatorname{Hom}_{\mathbf{A}}(Q,SN_i)$ (Corollary~1 to
Theorem~1, Chapter~II); the proof follows directly from these facts
(Lemma~2, \S2).

\begin{corollary}
Let $\mathbf{C}$ be a localizing subcategory of a locally finite
$\Univ$-category. Then $\mathbf{C}$ and $\mathbf{A}/\mathbf{C}$ are
locally finite $\Univ$-categories.
\end{corollary}

Suppose now that $\mathbf{A}$ is a locally noetherian
$\Univ$-category. Let $\mathbf{A}'$ be the full subcategory of
$\mathbf{A}$ defined by the noetherian objects of $\mathbf{A}$. If
$\mathbf{C}$ is a localizing subcategory of $\mathbf{A}$, we denote
by $\mathbf{C}'$ the full subcategory of $\mathbf{A}$ defined by the
noetherian objects of $\mathbf{C}$. We leave to the reader the task
of proving the following:

\begin{proposition}
The map $\mathbf{C}\mapsto\mathbf{C}'$ is a bijective map from the
set of localizing subcategories of $\mathbf{A}$ onto the set of thick
subcategories of $\mathbf{A}'$.
\end{proposition}

In particular, if $\mathbf{A}$ is a locally finite $\Univ$-category,
every localizing subcategory $\mathbf{C}$ is characterised by the
simple objects it contains. It may be said that
$\mathbf{A}/\mathbf{C}$ is obtained from $\mathbf{A}$ by neglecting a
certain number of simple objects.
